

%\fancyhf{} % clear all header and footer fields
%\fancyhead[RO,R]{\thepage} %RO=right odd, RE=right even
%\renewcommand{\headrulewidth}{0pt}

%\begin{center}
%    %\Large
%    %\textbf{Thesis Title}
%    %
%    %\vspace{0.4cm}
%    %
%    %\vspace{0.4cm}
%    %\textbf{Thesis Author}
%    
%    \vspace{0.9cm}
%    \textbf{Abstract}
%    
%\end{center}

The increasing integration of composite materials into marine lifting surfaces, such as hydrofoils and turbine blades, 
is unlocking new opportunities for weight reduction and structural efficiency. However, this trend amplifies the need for reliable prediction 
of unsteady hydrodynamic behaviour and coupled hydro-elastic instabilities, with flutter remaining one of the most critical challenges. 
Most available numerical tools are either accurate but computationally expensive or rely on oversimplified modelling that fails to capture the complex 
coupling between hydrodynamics, structural anisotropy, and dynamic stability.
This work introduces a comprehensive framework that directly addresses these limitations. 
The toolchain integrates four key components into a unified workflow: a finite-element pre-processor for generating fully coupled 
six degrees of freedom beam models; a dynamic beam solver for time-domain structural response; a Doublet Lattice Method implementation for unsteady 
hydrodynamic load prediction; and a robust P-K analysis module featuring automatic mode tracking. 
Together, these modules enable accurate coupling of structural and hydrodynamic phenomena, delivering a comprehensive toolset for frequency-domain response 
evaluation and instability prediction.
The proposed framework, using reduced-order modelling, achieves high computational efficiency while maintaining accuracy.
This balance between fidelity and cost-effectiveness makes the tool ideal for early-stage design, parametric studies, and optimisation tasks.
The methodology and software architecture are described in detail and validated against experimental data for a hydrofoil with heterogeneous cross-sections.
Results show very good agreement, confirming the tool’s capability to capture coupled fluid–structure dynamics and accurately predict critical flutter conditions.
